\section{Related work}\label{sec:related}
\paragraph{Diffusion sampling and guidance.} We use denoising diffusion models \citep{ho2020ddpm}, deterministic DDIM sampling \citep{song2021ddim}, latent diffusion \citep{rombach2022ldm} and classifier-free guidance \citep{ho2022cfg}. Restricting guidance to an interval of noise levels has been reported to improve sample quality \citep{kynkaanniemi2024interval}, annealing the conditioning signal early to restore diversity \citep{sadat2024cads}, and choosing per-step guidance online with latent evaluators to be prompt-dependent \citep{dynamiccfg2025}. Add-noise-then-denoise operations are used for editing and for sampling \citep{meng2022sdedit,xu2023restart}; we use re-noising only as a \emph{measurement}.

\paragraph{Stages of generation.} Switching prompts over denoising intervals localises layout, shape and detail to successive stages \citep{patashnik2023mixing}; stage-specific text embeddings control different attributes \citep{zhang2023prospect}; an ensemble of stage-specialised denoisers is motivated by text conditioning mattering most early \citep{balaji2022ediffi}; and features have been analysed theoretically as emerging in narrow time windows \citep{li2024critical}. Strengthening the conditioning signal at selected steps has been reported to improve alignment more when done late than early \citep{grimal2025step}, a finding in tension with front-loading guidance, and negative prompts to act with a delay \citep{ban2024negative}.

\paragraph{Cross-attention.} Cross-attention maps are used to interpret \citep{tang2023daam} and edit \citep{hertz2023p2p} generations, and attention-based objectives target object neglect and attribute binding \citep{chefer2023attend,li2023divide,fu2024gather,rassin2023syngen}; attention can also read and steer object size and position \citep{epstein2023selfguidance}. Initial-noise search and early prediction of image quality from attention are close in aim to our attention analysis \citep{guo2024initno,diffusionprobe2026}.

\paragraph{Evaluation.} Object-focused and question-answering evaluators \citep{ghosh2023geneval,huang2023t2icompbench,hu2023tifa,kamath2025geneval2} motivate our emphasis on judge reliability. Our judge combines open-vocabulary detection \citep{minderer2022owlvit} with CLIP \citep{radford2021clip}; seed diversity uses DINOv2 embeddings \citep{oquab2024dinov2}.

\paragraph{Position of this work.} Table~\ref{tab:related} contrasts what the closest methods measure, intervene on and evaluate. The common thread is that timing matters; the differences are that we (a) measure timing \emph{interventionally and per constraint type} with an explicit chance correction, (b) test whether CFG moves that timing, (c) compare an unsupervised attention signal against the interventional one, and (d) derive a guidance switch from the measurement rather than tuning it on adherence. To our knowledge these four elements have not been combined on controlled scene graphs; the literature search was narrow (Table~\ref{tab:related} caption) and this is a candidate gap, not an established one.
\begin{table*}[t]\centering\small
\caption{Closest prior work and this study: what each method measures, what it changes during sampling, and how it is evaluated. ``--'': not an aim of the method. Descriptions of Prompt-to-Prompt and Attend-and-Excite were checked against the full papers, the others against their abstracts (marked in the bibliography source).}\label{tab:related}
\begin{tabular}{p{2.9cm}p{3.7cm}p{3.7cm}p{3.7cm}}
\toprule
Method & Measures & Intervenes on & Evaluated by \\
\midrule
Prompt-to-Prompt \citep{hertz2023p2p} & -- & cross-attention maps (injection, re-weighting) to edit an image & edit fidelity, qualitative and quantitative \\
Attend-and-Excite \citep{chefer2023attend} & subject-token attention as a neglect signal & latents at inference, to raise neglected tokens' attention & object presence and binding (CLIP-based metrics, user study) \\
Prompt Mixing \citep{patashnik2023mixing} & -- & which prompt is used in which denoising interval & localisation of shape variations \\
Critical windows \citep{li2024critical} & time windows in which features emerge (theory) & -- & theory and experiments \\
Interval guidance \citep{kynkaanniemi2024interval} & -- & CFG applied only in a limited noise interval & distribution and sample quality \\
CADS \citep{sadat2024cads} & -- & conditioning annealed with noise early in sampling & diversity and quality \\
Dynamic CFG \citep{dynamiccfg2025} & per-step feedback from latent evaluators & per-step guidance scale, chosen online & prompt adherence and quality \\
Step selection \citep{grimal2025step} & -- & conditioning signal strengthened at selected steps & text--image alignment \\
InitNO \citep{guo2024initno} & attention-based validity of the initial noise & initial noise, optimised before sampling & text--image alignment \\
Diffusion Probe \citep{diffusionprobe2026} & early cross-attention as a predictor of final quality & -- (used for prompt and seed selection) & prediction accuracy \\
\midrule
This study & interventional, chance-corrected persistence $C(t)$ and $t_{50}$ per constraint type; attention stabilisation & two-level CFG switched at a step derived from development-set $t_{50}$ (CAG) & pre-registered tests with a human-audited judge; constant, reversed and matched-mean schedule controls \\
\bottomrule
\end{tabular}
\end{table*}

